% ===========================================================================
%  oscNext L4 -- noise BDT status
%
%  Compile:  pdflatex oscnext_l4_noise_bdt.tex   (twice, for the frame numbers)
%
%  Figures go in ./figures/ .  A figure that is not there yet is replaced by a
%  box of the SAME SIZE naming the file it wants, so the layout does not move
%  when you drop the real plot in, and the deck compiles at every stage.
%  See FIGURES.md for which plot belongs where.
%
%  STYLE RULES for this deck -- keep them when editing:
%    * plots carry the talk; text is a frame around them, never a paragraph
%    * NO numbers in running prose.  Every number lives in a tabular or in a
%      plot, where it can be read off and pointed at.
%    * anything about pass2 is sourced from the technical note and says so.
%      Nothing here reprocesses pass2.
% ===========================================================================
\documentclass[aspectratio=169,10pt]{beamer}

\usetheme{default}
\setbeamertemplate{navigation symbols}{}
\setbeamertemplate{footline}[frame number]
\setbeamercolor{frametitle}{fg=black}
\setbeamerfont{frametitle}{size=\large,series=\bfseries}
\setbeamerfont{framesubtitle}{size=\small}
\setbeamercolor{framesubtitle}{fg=gray}

\usepackage[utf8]{inputenc}   % harmless on TeX Live >=2018, needed before
\usepackage[T1]{fontenc}
\usepackage{graphicx}
\usepackage{booktabs}
\usepackage{array}        % >{\raggedright\arraybackslash}p{} columns
\usepackage{amsmath}
\usepackage{amssymb}      % \checkmark
\usepackage{xcolor}

\definecolor{good}{RGB}{0,110,55}
\definecolor{bad}{RGB}{170,30,30}
\definecolor{warn}{RGB}{185,105,0}

% --- figure with a fixed vertical budget ----------------------------------
% \plotht is the height every figure box gets.  Set it per frame BEFORE the
% columns; the image is scaled to fit inside it and the placeholder occupies
% exactly the same space, so a missing figure never shifts the layout.
\newlength{\plotht}
\setlength{\plotht}{0.42\textheight}

\newcommand{\plot}[1]{%
  \IfFileExists{figures/#1}%
    {\includegraphics[height=\plotht,width=\linewidth,keepaspectratio]{figures/#1}}%
    {\fbox{\parbox[c][\plotht][c]{0.93\linewidth}{%
        \centering\ttfamily\scriptsize \detokenize{figures/#1}\\[2pt]
        \normalfont\itshape\scriptsize(not present yet)}}}}

\newcommand{\figcap}[1]{\\[1pt]{\scriptsize #1}}
\newcommand{\tick}{\textcolor{good}{\textbf{\checkmark}}}
\newcommand{\cross}{\textcolor{bad}{\textbf{$\times$}}}
\newcommand{\src}[1]{{\scriptsize\itshape\color{gray}#1}}

\title{oscNext Level 4: rebuilding the noise classifier}
\subtitle{pass3 processing, variable validation, and the choice of BDT}
\author{Koray Gökçeler}
\date{\today}

\begin{document}
% =========================================== 1  TITLE ========================
\begin{frame}[plain]
  \titlepage
  \begin{center}\small
  L3 $\to$ L4 rebuilt for pass3 \quad$\cdot$\quad
  noise BDT trained \quad$\cdot$\quad
  the note's target reached
  \end{center}
\end{frame}

% =========================================== 2  LEVELS =======================
\begin{frame}[c]{The oscNext selection chain}
  \framesubtitle{what each level is for -- section 3 of the technical note}
  \centering\scriptsize
  \begin{tabular}{@{}>{\raggedright\arraybackslash}p{1.7cm}%
                    >{\raggedright\arraybackslash}p{11.9cm}@{}}
    \toprule
    \textbf{Level} & \textbf{Objective, as the note states it} \\
    \midrule
    \textbf{L2}\newline\scriptsize DeepCore Filter &
      The chain starts from the collaboration-wide Level 2.  For the low-energy
      working group ``the only filter which matters for the event selection is
      the DeepCore Filter''.  It splits the SRT-cleaned pulses into a DeepCore
      fiducial and a DeepCore veto series, and discards veto hits that could be
      causally related to a muon crossing the detector. \\[3pt]
    \textbf{L3} &
      ``The objective of Level3 is to use simple, fast cuts to remove regions
      of data-MC disagreement that result from event types for which we do not
      have MC'' -- muon bundles, coincident events, SLOP triggers.  Rewritten
      for oscNext to avoid any dependence on PMT charge. \\[3pt]
    \textbf{\textcolor{bad}{L4}} &
      ``After the Level 3 selection, the event rate is under 1\,Hz.  Crucially,
      the data-MC agreement of the sample is much improved, allowing machine
      learning methods to be applied for efficient background rejection.''
      \textbf{Two classifiers: noise and muon.} \\[3pt]
    \textbf{L5} &
      The main goal of Level 5 is ``to further reduce the remaining muon
      background, achieving a neutrino dominated sample'' -- starting
      containment cuts, and corridor cuts against muons that sneak between the
      IceCube strings. \\
    \bottomrule
  \end{tabular}\\[3pt]
  \src{quoted from sections 3.4--3.7; reconstruction and PID follow in
       section 4.  The note's selection chapter stops at L5.}

  \medskip
  \begin{columns}[c,onlytextwidth]
    \begin{column}{0.42\textwidth}
      \centering\scriptsize
      \begin{tabular}{@{}lr@{}}
        \toprule
        at L3 & rate [mHz] \\
        \midrule
        atmospheric muons & 505 \\
        detector noise    & 36.6 \\
        neutrinos         & 5.25 \\
        \bottomrule
      \end{tabular}\\[2pt]
      \src{technical note, Table 13}
    \end{column}
    \begin{column}{0.54\textwidth}
      \footnotesize
      \textbf{L4 inverts this ratio} -- and has to do it before the
      reconstruction, which only starts at L5.
    \end{column}
  \end{columns}
  \note{The design principle is in section 3.1: the first levels use simple
  and fast algorithms to bring the rate down to where more expensive ones can
  be applied.  That is why L4 may not reconstruct: reconstruction is what L5
  onward pays for.  Note there is no Level 6 or 7 in this document -- the
  selection chapter ends at L5 and reconstruction/PID is a separate chapter.}
\end{frame}

% =========================================== 3  L4 ===========================
\begin{frame}[c]{What L4 does, and what had to be rebuilt}
  \begin{columns}[T,onlytextwidth]
    \begin{column}{0.48\textwidth}
      \footnotesize
      \[
        \text{L4} \;=\; P_{\nu}^{\text{noise}} \ge 0.70
        \;\wedge\;
        P_{\nu}^{\text{muon}} \ge 0.65
      \]
      \centering\src{thresholds: technical note, sec.\ 3.5}

      \bigskip
      \raggedright
      \textbf{Why rebuild?}  The original segment is commented out and needs
      missing projects:
      \begin{itemize}\footnotesize
        \itemsep2pt
        \item \texttt{icecube.oscNext} -- \texttt{I3Classifier}
        \item \texttt{tau\_bdt} -- VICH
        \item \texttt{analysis.event\_selection} -- Dunkman variables
      \end{itemize}

      \medskip
      Part of the inputs were \textbf{rewritten in pure Python}, from the note.
    \end{column}
    \begin{column}{0.50\textwidth}
      \centering\footnotesize
      \textbf{Where the BDT inputs come from}\\[5pt]
      \begin{tabular}{@{}cl>{\raggedright\arraybackslash}p{2.9cm}@{}}
        \toprule
        n & source & status \\
        \midrule
        4 & ready from L3    & validated \\
        5 & IceTray projects & validated \\
        \textbf{5} & \textbf{our rewrite} & \textcolor{bad}{unverified} \\
        \bottomrule
      \end{tabular}\\[2pt]
      \src{definitions: note, Tables 11--12}

      \bigskip
      \raggedright
      \textbf{All the risk is in the last row.}

      \medskip
      Compared line by line against the note and the original pass2 segment --
      which found bugs that \textbf{raise no error}.
    \end{column}
  \end{columns}
  \note{Do not enumerate the bugs.  If asked: the index table overwriting
  data, non-unique event IDs mismatching rows, the flux event count frozen
  across files, and a bypassed cleaning step in micro\_count -- that last one
  is inherited from the original pass2 code, not introduced by us.}
\end{frame}

% =========================================== 4  VARIABLES ====================
\begin{frame}[c]{Noise BDT inputs, as the note defines them}
  \framesubtitle{Table 11 of the technical note, quoted}
  \centering\scriptsize
  \begin{tabular}{@{}>{\ttfamily\raggedright\arraybackslash}p{4.8cm}%
                    >{\raggedright\arraybackslash}p{8.7cm}@{}}
    \toprule
    \normalfont\textbf{Variable} & \textbf{Description} \\
    \midrule
    IC2018\_LE\_L3\_Vars.NchCleaned &
      ``See description in section 3.5.1'' --- there:
      ``A number of hit DOMs in the cleaned pulse series.'' \\[4pt]
    L4\_micro\_count.\-STW\_m3500p4000\_DTW200 &
      ``Start with the cleaned pulse series.  Look at pulses occurring within
      [$-3.5\,\mu$s, $+4\,\mu$s] from the trigger time.  Slide a time window
      of 200\,ns that maximizes the number of triggered DOMs in it.  Get the
      number of triggered DOMs in that sliding time window'' \\[4pt]
    L4\_iLineFit.speed &
      ``Speed fitted by the improved LineFit algorithm'' \\[4pt]
    L4\_fill\_ratio.\-fill\_ratio\_from\_mean &
      ``Measure of the geometrical spread of the hits about some vertex
      \textcolor{bad}{(details here)}'' \\[4pt]
    IC2018\_LE\_L3\_Vars.\-FullTimeLengthRatio &
      ``Measure the total duration of an event in both the cleaned and
      uncleaned pulse series (duration $=$ max pulse time $-$ min pulse
      time).  Calculate the ratio between the two.'' \\
    \bottomrule
  \end{tabular}\\[3pt]
  \src{descriptions quoted verbatim from Table 11; the NchCleaned line points
       at sec.\ 3.5.1, quoted beside it}

  \medskip
  \footnotesize\raggedright
  Two of these are \textbf{not definitions}: \texttt{fill\_ratio}'s vertex is
  an unfilled cross-reference, and the ratio's direction is never stated.
  Both were recovered from the original pass2 code and from Figure 13.
  \note{The red "(details here)" is in the note itself -- an unfilled
  cross-reference.  It matters because fill\_ratio turns out to carry most of
  the model's gain, so the one variable we lean on hardest is the one the note
  does not define.  We took the vertex from the original code: the position of
  the first HLC hit.  The ratio's direction came from the x axis of Figure 13,
  which runs 0 to 1 -- so it is cleaned over uncleaned.}
\end{frame}

% =========================================== 5  INPUT PLOT ===================
\setlength{\plotht}{0.80\textheight}
\begin{frame}[c]{Noise BDT inputs, signal vs.\ noise}
  \framesubtitle{pass3, our own processing}
  \begin{center}
    \plot{4.png}
    \figcap{rate vs.\ variable for the five inputs}
  \end{center}

  \vspace{2pt}
  {\footnotesize
   \textbf{Measured, not assumed.}  \texttt{FullTimeLengthRatio} separates the
   \emph{opposite} way to the expectation.}
  \note{The docstring used to claim the ratio is near one for a real event and
  near zero for noise.  Wrong: the uncleaned series spans the whole readout
  window in every event, so the variable is effectively the cleaned duration
  over a constant.  It does separate -- but because a low-energy cascade is
  compact in time, not because noise is.  If asked for the comparison with the
  note: this is the plot to hold beside Figure 13.}
\end{frame}

% =========================================== 6  CORRELATION ==================
\setlength{\plotht}{0.74\textheight}
\begin{frame}[c]{Input correlations}
  \framesubtitle{does any pair of the five inputs say the same thing?}
  \begin{columns}[c,onlytextwidth]
    \begin{column}{0.49\textwidth}
      \centering
      \plot{input_correlation_signal.png}
      \figcap{Spearman rank correlation, signal}
    \end{column}
    \begin{column}{0.49\textwidth}
      \centering
      \plot{input_correlation_background.png}
      \figcap{Spearman rank correlation, background}
    \end{column}
  \end{columns}
  \note{Two readings, both off the plots.  Left: the inputs are largely
  independent, so each one carries its own information and none is redundant.
  Right: fill_ratio carries about sixty percent of the gain and
  iLineFit_speed about one percent -- nearly dead.  That is the argument for
  next step 2: the variable the model leans on hardest is the one whose only
  free parameter, the spherical radius, is inherited from GRECO.  The
  original code says so itself -- "was optimised for GRECO but has not been
  re-optimised for oscNext".  Re-tuning it is the single highest-value knob
  we have.  If asked about the ordering: the importance comes from the
  trained model, so strictly it is a result of the next two slides; it sits
  here because it belongs with the inputs it is about.}
\end{frame}

% =========================================== 7  IMPORTANCE ===================
\setlength{\plotht}{0.78\textheight}
\begin{frame}[c]{Feature importance}
  \framesubtitle{which input does the work?}
  % Nudged left of centre: the long tick labels hang off the left of the bars,
  % so a centred figure reads as sitting too far right.  Centred would be an
  % indent of 0.06\textwidth; 0.03 puts it half that.
  \noindent\makebox[\textwidth][l]{%
    \hspace*{0.03\textwidth}%
    \begin{minipage}{0.88\textwidth}
      \centering
      \plot{feature_importance.png}
      \figcap{gain split of the trained noise model}
    \end{minipage}}
  \note{This is the argument for next step 2.  fill_ratio carries about sixty
  percent of the gain and iLineFit_speed about one percent -- nearly dead.
  The variable the model leans on hardest is the one whose only free
  parameter, the spherical radius, is inherited from GRECO: the original code
  says so itself, "was optimised for GRECO but has not been re-optimised for
  oscNext".  Re-tuning it is the single highest-value knob we have.  If asked
  about the ordering: the importance comes from the trained model, so strictly
  it is a result of the next two slides; it sits here because it belongs with
  the inputs it is about.}
\end{frame}

% =========================================== 8  BDT ==========================
\setlength{\plotht}{0.56\textheight}
\begin{frame}[c]{How a boosted decision tree works}
  \begin{columns}[T,onlytextwidth]
    \begin{column}{0.42\textwidth}
      \footnotesize
      Many shallow trees in sequence, each correcting the last.

      \medskip
      The engines differ in \emph{how} the next tree is chosen:

      \smallskip
      \centering
      \begin{tabular}{@{}>{\raggedright\arraybackslash}p{1.9cm}%
                        >{\raggedright\arraybackslash}p{3.2cm}@{}}
        \toprule
        \textbf{AdaBoost}\newline(pybdt) &
        re-weights the \emph{events} it gets wrong \\[4pt]
        \textbf{Gradient}\newline\textbf{boosting}\newline(LightGBM) &
        fits the \emph{gradient} of the loss \\
        \bottomrule
      \end{tabular}\\[2pt]
      \src{the note uses LightGBM, sec.\ 3.6.1}
    \end{column}
    \begin{column}{0.56\textwidth}
      \centering
      \plot{5.png}
      \figcap{trees in sequence, leaf values summed}

      \bigskip
      \raggedright\footnotesize
      \textbf{Output scale.}  LightGBM returns a probability, so the note's
      cut transfers to our model directly.  A pybdt score does not.
    \end{column}
  \end{columns}
  \note{Keep this short, the audience knows what a BDT is.  The one point that
  matters downstream is the output scale -- it decides whether the note's cut
  value can be reused or has to be recalibrated.}
\end{frame}

% =========================================== 9  PYBDT ========================
\setlength{\plotht}{0.42\textheight}
\begin{frame}[c]{Training with pybdt (AdaBoost)}
  \framesubtitle{a deliberate deviation from the note -- and its cost}
  % This figure is about five times wider than it is tall.  In a side column
  % it collapsed into an unreadable strip, so it gets the full frame width
  % and the table sits underneath it.
  \begin{center}
    \plot{2.png}
    \figcap{six configurations, efficiency vs.\ rejection, test set}
  \end{center}

  \vspace{2pt}
  \begin{columns}[T,onlytextwidth]
    \begin{column}{0.44\textwidth}
      \centering\scriptsize
      \begin{tabular}{@{}lrrr@{}}
        \toprule
        model & 90\,\% & 95\,\% & 99\,\% \\
        \midrule
        stump (d1)      & 91.3 & 76.0 & 76.0$^*$ \\
        \textbf{d2t500} & \textbf{94.0} & \textbf{91.7} & 62.0 \\
        d3t300          & 88.9 & 75.2 & 43.7 \\
        d4t500          & 91.5 & 86.1 & \textbf{65.8} \\
        d6t500          & 13.4 & 10.3 &  5.8 \\
        \bottomrule
      \end{tabular}\\[2pt]
      \src{signal efficiency [\%] at fixed noise rejection;
           $^*$too few distinct scores at depth one}
    \end{column}
    \begin{column}{0.54\textwidth}
      \footnotesize
      \cross~\textbf{the note's target is out of reach}

      \smallskip
      \tick~training is deterministic

      \smallskip
      \cross~the KS $p$-value is useless here -- we use the train/test
      efficiency \emph{gap}
    \end{column}
  \end{columns}
  \note{Be honest about the framing: these hyperparameters were hand-picked,
  not tuned to the same effort as the note's.  So this is not "AdaBoost is a
  worse algorithm" -- it is "our pybdt setup costs us a lot".  On the KS test:
  it flagged the best model and cleared the worst one.}
\end{frame}

% =========================================== 10 LIGHTGBM =====================
\setlength{\plotht}{0.70\textheight}
\begin{frame}[c]{Training with LightGBM}
  \framesubtitle{the note's own method and hyperparameters (Table 10)}
  \begin{columns}[T,onlytextwidth]
    \begin{column}{0.58\textwidth}
      \centering
      \plot{3.png}
      \figcap{same events, same variables, same split}
    \end{column}
    \begin{column}{0.40\textwidth}
      \centering\scriptsize
      \begin{tabular}{@{}lrr@{}}
        \toprule
        rejection & AdaBoost & \textbf{LightGBM} \\
        \midrule
        90\,\% & 94.0 & \textbf{99.0} \\
        95\,\% & 91.7 & \textbf{98.5} \\
        99\,\% & 65.8 & \textbf{95.9} \\
        \bottomrule
      \end{tabular}\\[2pt]
      \src{signal efficiency [\%], test set, event counts}

      \bigskip
      \raggedright\footnotesize
      \tick~\textbf{the note's target is reached}

      \medskip
      \tick~no overtraining -- train and test coincide

      \medskip
      \tick~reached while \emph{handicapped} by the note's leaf-size floor
    \end{column}
  \end{columns}
  \note{The handicap is worth making explicit: the note's leaf-size floor was
  tuned on far larger noise statistics than ours, and training stopped early
  well short of the tree budget.  The model met the target anyway.}
\end{frame}

% =========================================== 11 SCORE DIST ===================
\setlength{\plotht}{0.78\textheight}
\begin{frame}[c]{LightGBM score distribution}
  \framesubtitle{the overtraining check}
  \begin{center}
    \plot{lightgbm_score_dist.png}
    \figcap{train and test lie on top of each other}
  \end{center}
  \note{Left panel linear, right panel the same histogram on a log scale --
  the tail between the two peaks is invisible without it, and that tail is
  where the overtraining would show.  Train and test coincide in both.  The
  y axis is the training weight summed per bin: each class is normalised to
  the same total, so the two areas are equal by construction -- do not read
  it as "there is as much background as signal".  For noise the weights are
  uniform, so the background curves are effectively event counts.  Background
  stops around 0.85: no background event scores higher than that.}
\end{frame}

% =========================================== 12 MC LIMIT =====================
\setlength{\plotht}{0.78\textheight}
\begin{frame}[c]{The limit: noise MC statistics}
  \framesubtitle{why the far tail is not a measurement}
  \begin{center}
    % swap for mc_scaling.png if the background-scaling scan gets produced
    \plot{lightgbm_cuts.png}
    \figcap{grey lines: where the background statistics runs out}
  \end{center}
  \note{The numbers, if anyone asks.  The test set holds 329,545 signal events
  against 2,051 noise -- a ratio of 159 to 1.  Above the cut there are about
  100 background events left at 90 percent rejection, 50 at 95, ten at 99,
  five at 99.5 and one at 99.9.  The note's target sits at 99.2 percent
  rejection, which our test set defines with eight events.  So: not the
  bottleneck it looked like -- an earlier conclusion held that background
  statistics limited the model and the LightGBM measurement retracted it.
  What we have is enough to REACH the target, not enough to MEASURE the far
  tail.  More vuvuzela MC is valuable but not blocking, and stating it the
  other way round would overstate the request.}
\end{frame}

% =========================================== 13 TO DO ========================
\begin{frame}[c]{Status and next steps}
  \begin{columns}[T,onlytextwidth]
    \begin{column}{0.44\textwidth}
      \footnotesize
      \textbf{Done}
      \begin{itemize}\footnotesize
        \item L3 $\to$ L4 rebuilt for pass3, robust against corrupt files
        \item all BDT inputs located
        \item validated against the note and the original segment
        \item noise classifier trained, target reached
      \end{itemize}
    \end{column}
    \begin{column}{0.54\textwidth}
      \footnotesize
      \textbf{Next}
      \begin{enumerate}\footnotesize
        \item \textbf{Cross-check against pass2 L4 files} -- validates the
              rewritten variables; tests each model on the other pass
        \item \textbf{Re-tune the \texttt{fill\_ratio} radius} -- carries most
              of the gain (see the importance plot), never re-optimised for
              oscNext
        \item \textbf{Muon classifier} -- CORSIKA background
        \item \textbf{More vuvuzela MC} -- to measure the far tail
        \item Reprocess the remaining samples
      \end{enumerate}
    \end{column}
  \end{columns}

  \vfill
  {\scriptsize
   \textbf{Limitations:} no $\nu_\tau$ set; the muon background is CORSIKA,
   so no data/MC check yet.}
  \note{End on item 1 -- the pass2 cross-check closes the remaining validation
  gap, and it is the natural ask if anyone can point at where those files
  live.  Item 2 comes straight off the plot on this slide.}
\end{frame}

\end{document}
